
\documentclass{article}
\usepackage{colortbl}
\usepackage{makecell}
\usepackage{multirow}
\usepackage{supertabular}

\begin{document}

\newcounter{utterance}

\centering \large Interaction Transcript for game `cladder', experiment `full\_v1.5\_default', episode 855 with Qwen3.6{-}27B.
\vspace{24pt}

{ \footnotesize  \setcounter{utterance}{1}
\setlength{\tabcolsep}{0pt}
\begin{supertabular}{c@{$\;$}|p{.15\linewidth}@{}p{.15\linewidth}p{.15\linewidth}p{.15\linewidth}p{.15\linewidth}p{.15\linewidth}}
    \# & \multicolumn{2}{c}{Player} && \multicolumn{2}{c}{Game Master} \\
    \hline

    \theutterance \stepcounter{utterance}  
    & & & \multicolumn{4}{p{0.6\linewidth}}{
        \cellcolor[rgb]{0.9,0.9,0.9}{
            \makecell[{{p{\linewidth}}}]{
                \texttt{\tiny{[P1$\langle$GM]}}
                \texttt{You are an expert in causal inference. The following question is not a typical commonsense query, but rather a meticulously designed question created by a professor specializing in causal inference, intended to assess the students' mastery of the course content.} \\
\texttt{Imagine a self{-}contained, hypothetical world with only the following conditions, and without any unmentioned factors or causal relationships: Personality has a direct effect on relationship status. Appearance has a direct effect on relationship status. The overall probability of kindness is 85\%. For people with mean personalities and single people, the probability of attractive appearance is 38\%. For people with mean personalities and in a relationship, the probability of attractive appearance is 71\%. For people with kind personalities and single people, the probability of attractive appearance is 32\%. For people with kind personalities and in a relationship, the probability of attractive appearance is 57\%. If we look at people in a relationship, does the chance of attractive appearance decrease when kindness?} \\
\texttt{Start your answer with 'Yes' or 'No', followed by additional reasoning or evidence to support your explanation.} \\
            }
        }
    }
     \\ \\

    \theutterance \stepcounter{utterance}  
    & \multicolumn{4}{p{0.6\linewidth}}{
        \cellcolor[rgb]{0.561,0.737,0.561}{
            \makecell[{{p{\linewidth}}}]{
                \texttt{\tiny{[P1$\rangle$GM]}}
                \texttt{No} \\
\\ 
\texttt{**Reasoning:**} \\
\\ 
\texttt{To determine whether the chance of attractive appearance decreases when kindness is present among people in a relationship, we must compare the conditional probability of attractive appearance given that the person is in a relationship and kind, versus the probability given that the person is in a relationship and mean.} \\
\\ 
\texttt{From the problem description, we are given the following specific conditional probabilities for people **in a relationship**:} \\
\\ 
\texttt{1.  For people with **mean** personalities and **in a relationship**, the probability of attractive appearance is **71\%** (\$P(\textbackslash{}text\{Attractive\} | \textbackslash{}text\{Mean\}, \textbackslash{}text\{Relationship\}) = 0.71\$).} \\
\texttt{2.  For people with **kind** personalities and **in a relationship**, the probability of attractive appearance is **57\%** (\$P(\textbackslash{}text\{Attractive\} | \textbackslash{}text\{Kind\}, \textbackslash{}text\{Relationship\}) = 0.57\$).} \\
\\ 
\texttt{The question asks: "If we look at people in a relationship, does the chance of attractive appearance decrease when kindness?"} \\
\\ 
\texttt{This phrasing implies a comparison between the state of being kind and the alternative state (being mean) within the subset of people who are in a relationship.} \\
\\ 
\texttt{*   Probability of attractive appearance for **Kind** people in a relationship: **57\%**} \\
\texttt{*   Probability of attractive appearance for **Mean** people in a relationship: **71\%**} \\
\\ 
\texttt{Comparing the two: \$57\textbackslash{}\% < 7} \\
            }
        }
    }
    & & \\ \\

    \theutterance \stepcounter{utterance}  
    & & & \multicolumn{2}{p{0.3\linewidth}}{
        \cellcolor[rgb]{0.9,0.9,0.9}{
            \makecell[{{p{\linewidth}}}]{
                \texttt{\tiny{[GM$|$GM]}}
                \texttt{No} \\
            }
        }
    }
    & & \\ \\

    \theutterance \stepcounter{utterance}  
    & & & \multicolumn{2}{p{0.3\linewidth}}{
        \cellcolor[rgb]{0.9,0.9,0.9}{
            \makecell[{{p{\linewidth}}}]{
                \texttt{\tiny{[GM$|$GM]}}
                \texttt{yes} \\
            }
        }
    }
    & & \\ \\

    \theutterance \stepcounter{utterance}  
    & & & \multicolumn{2}{p{0.3\linewidth}}{
        \cellcolor[rgb]{0.9,0.9,0.9}{
            \makecell[{{p{\linewidth}}}]{
                \texttt{\tiny{[GM$|$GM]}}
                \texttt{game\_result = LOSE} \\
            }
        }
    }
    & & \\ \\

\end{supertabular}
}

\end{document}
